\documentclass[10pt,a4paper]{amsart}
\usepackage[margin=19mm]{geometry}
\usepackage{amsmath,amssymb,mathtools}
\usepackage[scr=rsfs,cal=euler]{mathalfa}
\usepackage{xfrac}
\usepackage[unicode,hidelinks,bookmarks=false]{hyperref}
\newcommand{\ie}{i.e.\ }
\newcommand{\cf}{cf.\ }
\newcommand{\resp}{respectively\ }
\newcommand{\Subsection}{\subsection}
\providecommand{\nobreakdash}{\nobreak\hbox{-}}
\setlength{\emergencystretch}{2em}
\multlinegap0pt
\usepackage{bm}
\usepackage{bbm}
\usepackage{dsfont}
\usepackage{xspace}
\usepackage{cancel}
\usepackage{mathtools}
\usepackage{amsthm}
\makeatletter
\@ifundefined{defi}{\newtheorem{defi}{Defi}}{}
\@ifundefined{prop}{\newtheorem{prop}[defi]{Proposition}}{}
\makeatother
\newcommand{\BO}{\mathcal{O}}
\newcommand{\C}{\mathbb{C}}
\newcommand{\Z}{\mathbb{Z}}
\newcommand{\p}{\varphi}
\newcommand{\sz}[1]{\left| \vec{#1} \right|}
\title[Rational approximation of Euler's constant using multiple or]{Rational approximation of Euler's constant using multiple orthogonal polynomials}
\author{Thomas Wolfs, Walter Van Assche}
\date{Source: arXiv:2404.09799v3 (CC BY 4.0)}
\begin{document}
\maketitle

\noindent\emph{Source and licence.} Rational approximation of Euler's constant using multiple orthogonal polynomials. Version arXiv:2404.09799v3 (CC BY 4.0), distributed under the Creative Commons Attribution 4.0 International licence, as recorded in the arXiv licence metadata for this exact version and shown on the abstract page https://arxiv.org/abs/2404.09799v3. Full rights notice: licenses/arxiv-2404.09799-CC-BY-4.0.txt.\par\smallskip

\noindent\textit{Editorial scope.} Counted unit: the Prop whose source label is \texttt{I\_AQ\_error} in the article (source0.tex lines 443--446, source label \texttt{I\_AQ\_error}), delivered together with the article's own complete original proof of it (source0.tex lines 447--469, one \texttt{proof} environment of 23 lines). No mathematics is written, repaired or re-derived: statement and proof are the article's own text line for line. Required context retained uncounted: I\_int (lines 303--307); I\_AQ\_BT (lines 402--405). Every definition, display and label of those lines is retained unchanged. Omitted from this package and not redistributed: 1--302, 308--401, 406--442, 470--831. Nothing inside a retained excerpt is omitted or altered: every retained body is delivered exactly as the article's lines are written, so no omission receipt of any kind accompanies this package. Citations: the retained excerpts carry no citation command at all, so no bibliography entry of the article is transcribed and no reference is redistributed. Reference adaptation: none was needed and none was made; every reference inside the delivered unit resolves inside it, so no unresolved reference or citation is produced. Printed slips kept literal: no printed word, symbol, hypothesis or number of the counted statement or of its proof was altered. Layout and class: the article's own class is \texttt{article} with packages including inputenc, babel, amsmath, amsthm, amssymb, mathtools, tikz, verbatim, fancyvrb, array, adjustbox, changepage, subcaption, breqn; the wrapper is the lane's pinned \texttt{amsart} editorial wrapper at 10pt on A4, which restates the theorem-like environments the retained text needs and adds the article's own macro definitions that the retained text needs (\texttt{\textbackslash BO}, \texttt{\textbackslash C}, \texttt{\textbackslash Z}, \texttt{\textbackslash p}, \texttt{\textbackslash sz}), with extra wrapper packages (bm, bbm, dsfont, xspace, cancel, mathtools). The delivered numbering is the unit's own. Assets: no retained span contains a figure environment, a graphics-inclusion command, a PDF-page inclusion, a table or a TikZ picture; no image file is copied into this package, the package declares no asset dependency and no image file is redistributed with it. Licence: version arXiv:2404.09799v3 of this article is distributed under the Creative Commons Attribution 4.0 International licence, as recorded in the arXiv licence metadata for this exact version and shown on the abstract page https://arxiv.org/abs/2404.09799v3. A full rights notice is delivered at licenses/arxiv-2404.09799-CC-BY-4.0.txt. Characters: every retained excerpt is pure ASCII, so the delivered engine, XeLaTeX with its default Latin Modern text font, reports no missing character.
\begin{prop} \label{I_int}
Suppose that $\sz{n}=\sz{m}+1$ and $m_1+1\geq m_2$. Assume that $\alpha-\beta\not\in\Z$. Then,
   $$ F_{\vec{n},\vec{m}}^{(I\mid \alpha,\beta,\nu)}(x) =  \frac{(-1)^{\sz{n}+1}}{m_2!} \int_\Sigma  \frac{(t+\nu+1)_{m_2} }{(\alpha-t)_{n_1}(\beta-t)_{n_2}} \frac{x^t}{\Gamma(t+1)} \frac{dt}{2\pi i}, $$
   in terms of a counterclockwise contour $\Sigma$ in $\{t\in\C\mid -1< \text{Re}(t) < \min\{\alpha,\beta\}\}$ that encloses $(0,\infty)$ exactly once.
\end{prop}

\begin{prop} \label{I_AQ_BT}
    There are four solutions $\{S_n^{(I)}(\omega) \mid \omega^4=x\}$ that form a basis over $\C$ for the space of solutions of the recurrence relation for $F_{n;2}^{(II)}(x)$ with asymptotic behavior
    $$S_{n}^{(I)}(\omega) = n^{-9/8} \exp\left(4\omega n^{3/4} - \frac{1}{2}\omega^2 n^{1/2} - \frac{3}{8}\omega^3 n^{1/4} \right) \left(1+\BO(n^{-1/4})\right),\quad n\to\infty.$$
\end{prop}

\begin{prop} \label{I_AQ_error}
    Suppose that $x>0$. Then,
    $$F_{n}^{(I)}(x) \in\text{span}\{\text{Re}(S_{n}^{(I)}(i\sqrt[4]{x})),\text{Im}(S_{n}^{(I)}(i\sqrt[4]{x})),S_n^{(I)}(-\sqrt[4]{x})\}.$$     
\end{prop}

\begin{proof}
By making use of the relation $(-t)_{n+1} = (-1)^{n+1} (t-n)_{n+1}$, we may write the contour integral in Proposition \ref{I_int} as 
    $$F_n^{(I)}(x) = n! \int_{\Sigma} f_n(t) x^t \frac{dt}{2\pi i},\quad f_n(t) = \frac{\Gamma(t+n+1)\Gamma(t-n)^2}{\Gamma(t+1)^4}.$$
Afterwards, we perform a change of variables of the form $t\mapsto n^{\alpha} s$ for some $0<\alpha<1$ that will be chosen later. We will now approximate the gamma functions in the integrand via Stirling's formula
    $$\Gamma(s+a) \sim (2\pi)^{\frac{1}{2}} e^{-s} s^{s+a-\frac{1}{2}},\quad |s|\to\infty,\ \arg(s) < \pi. $$
The notation $f(t)\sim g(t)$ as $t\to\infty$ means that $\lim_{t\to\infty }f(t)/g(t) = 1$. We find that
$$f_n(n^{\alpha}s) \sim \frac{1}{(2\pi)^{\frac{1}{2}}} e^{n^{\alpha}s+n} \frac{(n^{\alpha}s+n)^{n^{\alpha}s+n+\frac{1}{2}} (n^{\alpha}s-n)^{2n^{\alpha}s-2n-1}}{(n^{\alpha}s)^{4n^{\alpha}s+2}},\quad n\to\infty. $$
By collecting some factors of powers of $n$ in the ratio, we may rewrite it as 
$$ n^{(3-4\alpha)n^{\alpha}s - n - \frac{1}{2} -2\alpha} \frac{(1+n^{\alpha-1}s)^{n^{\alpha}s+n+\frac{1}{2}} (n^{\alpha-1}s-1)^{2n^{\alpha}s-2n-1}}{s^{4n^{\alpha}s+2}}. $$
Since in what follows, we would like to apply the saddle point method, it makes sense to take $\alpha=3/4$ so that the $n^\alpha \ln n$ term in the exponent disappears. Doing so yields,
$$f_n(n^{3/4}s) \sim \frac{e^{n}}{(2\pi)^{\frac{1}{2}} n^{n+2}} e^{n^{3/4}s} \frac{(1+n^{-1/4}s)^{n^{3/4}s+n+\frac{1}{2}} (n^{-1/4}s-1)^{2n^{3/4}s-2n-1}}{s^{4n^{3/4}s+2}}, $$
which we may simplify to
$$f_n(n^{3/4}s;x) \sim - \frac{n^{-3/2}}{n!} e^{n^{3/4}s} \frac{(1+n^{-1/4}s)^{n^{3/4}s+n} (1-n^{-1/4}s)^{2n^{3/4}s-2n}}{s^{4n^{3/4}s+2}}, $$
by applying Stirling's formula to $n!$ and by making use of the fact that $(1+n^{-1/4}s)^{\frac{1}{2}}\sim 1$ and $(n^{-1/4}s-1)^{-1}\sim -1$ as $n\to\infty$. Therefore,
$$ n! n^{3/4} f_n(n^{3/4}s) x^{n^{3/4}s} \sim - s^{-2} n^{-3/4} e^{n^{3/4}\p_n(s)},  $$
where
    $$\p_n(s) = s(1-4\ln s+\ln x) + (s+n^{1/4})\ln(1+n^{-1/4}s) + 2 (s-n^{1/4}) \ln(1-n^{-1/4}s).$$  
We can use Maple to look for approximate saddle points $s^\ast_n = a+bn^{-1/4}+cn^{-1/2}+O(n^{-3/4})$ such that $\p_n'(s^\ast_n)=O(n^{-3/4})$ as $n\to\infty$ (to this end, we will expand the logarithms). If we would only consider leading terms, we would have to look for the actual saddle points of $\p(s)=s(4-4\ln s+\ln x)$, which are given by the solutions of $s^4=x$. In general, we get a system of equations that we can solve for the values of the parameters. We then find four of such approximate saddle points
    $$s^\ast_n(\omega) = \omega - \frac{1}{4} \omega^2 n^{-1/4} - \frac{9}{32} \omega^3 n^{-1/2} + O(n^{-3/4}),\quad \omega^4=x, $$
and we have
$$\p_n(s^\ast_n(\omega)) = 4\omega -\frac{1}{2} \omega^2 n^{-1/4} - \frac{3}{8} \omega^3 n^{-1/2} + O(n^{-3/4}). $$
By applying the (modified) saddle point method, we can then show that $F_n^{(I)}(x)$ is asymptotic to an appropriate linear combination of $n^{-9/8}\exp\left(n^{3/4}\p_n(s^\ast_n(\omega))\right)$ for $\omega^4=x$. Observe that these possible asymptotics are exactly those found in Theorem \ref{I_AQ_BT} via the Birkhoff-Trjitzinsky theory. The specific asymptotics that appear in the combination here depend on the approximate saddle points that the curve of steepest descent, which the contour $n^{-3/4}\Sigma$ needs to be deformed to, passes through. Note that the contour $n^{-3/4}\Sigma$ can't be deformed to a curve that passes through the approximate saddle point $s_n^\ast(x^{1/4})$ on the positive real line because then it would need to cross the strip $[0,n^{1/4}]$ inside $n^{-3/4}\Sigma$ where $f_n(t)$ has its poles. Hence we must have the stated result. 
\end{proof}

\end{document}
